\textbf{13.} Докажите теорему о пространственной иерархии: если функции $f$ и $g$ конструируемы по памяти, $f(n) \ge \log{n}$ и $f(n) = o(g(n))$, то
\[
\mathbf{DSPACE}(f(n)) \ne \mathbf{DSPACE}(g(n)).
\]

\textbf{Решение.} Пусть $h(n) = \frac{f(n) + g(n)}{2}$. Тогда $h(n)$ конструируема по памяти за $O(g(n))$ и при этом $f(n) = o(h(n))$ и $h(n) = O(g(n))$. Пусть $M_1, M_2, \ldots$ --- главная нумерация всех машин Тьюринга. Рассмотрим машину $D$, которая на входе $1^n$ запускает $M_n(1^n)$, ограничивая её зону работы $h(n)$ ячейками, а на входе иного вида возвращает 0. В отличие от теоремы для времени, нужно отдельно проконтролировать, что машина не зациклится. Для этого нужно завести счётчик шагов и проверять, что он не превысил количества всех конфигураций --- $a^{h(n)} h(n) n q$, где $a$ -- размер алфавита, $q$ -- число состояний $M_n$, множитель $nh(n)$ отвечает за расположение двух указателей на двух лентах --- $n$ вариантов для входной ленты и $h(n)$ для рабочей. Если вычисление $M_n(1^n)$ закончилось с некоторым ответом, то $D$ возвращает противоположный. Если же вычисление зациклилось (т.е. счётчик был превышен), то $D$ возвращает любой ответ, например 0 для определённости.

Машина $D$ распознаёт некоторый язык $L$. Покажем, что
\[
L \in \mathbf{DSPACE}(g(n)) \setminus \mathbf{DSPACE}(f(n)).
\]
Вначале посчитаем, сколько места заняло вычисление. Зона работы машины $M_n$ искусственно ограничена $h(n)$ ячейками, счётчик дополнительно займёт $O(h(n))$ ячеек (здесь использовали $f(n) \ge \log{n}$, так как в максимальном значении есть множитель $n$, уже дающий $\log{n}$ знаков), а в силу определения $h$ вычисление максимального значения счётчика будет выполнено за $O(g(n))$ памяти. Моделирование на универсальной машине Тьюринга увеличит память в константу раз. Поскольку $h(n) = O(g(n))$, язык будет распознан на памяти $O(g(n))$. С другой стороны, пусть он распознаётся на памяти $Cf(n)$, где $C$ --- константа. Это делает некоторая машина $M_k$. Можно считать, что $k$ достаточно большое, чтобы выполнялось $Cf(k) < h(k)$ (напомним, в главной нумерации одинаковые машины Тьюринга встречаются бесконечно часто). Но в таком случае $D(1^k)$ вычисляет $M_k(1^k)$ и обращает результат. Поэтому $M_k$ не может распознавать тот же язык, что и $D$. Полученное противоречие показывает, что этот язык не лежит в $\mathbf{DSPACE}(f(n))$, что и требовалось.